\section{The System}
\label{sec:system}

\subsection{Build layer, run once per city on a free GPU}
For any area of interest, Google Earth Engine exports four free global datasets. FABDEM
\cite{hawker2022fabdem} gives ground elevation at 30\,m with forests and buildings stripped out, so
that the surface is what water actually runs over. ESA WorldCover \cite{zanaga2021worldcover} gives
land cover at 10\,m, which tells us whether a cell is concrete, soil or vegetation. JRC surface
water \cite{pekel2016} tells us which cells are permanently wet, so that a lake is not mistaken for
a flood. Sentinel-1 radar scenes \cite{torres2012} give observed water on real storm dates, and
radar is used because it sees through cloud, which matters when the thing being observed is a
monsoon.

The model domain is a square grid at 60\,m. Patna-family and Bengaluru bundles use $128^2$ cells,
which is a window about 7.7\,km across, and each Mumbai tile uses $256^2$ cells, about 15.4\,km.
The Mumbai tiles share exact latitude and longitude edges, so adding them up never double-counts
any ground. Those tiles are independent models with closed boundaries, which means water does not
flow from one tile into its neighbour and tides are not simulated. The dashboard says so on the
page.

The solver is the local inertial shallow water scheme of Bates \emph{et al.} \cite{bates2010},
written in PyTorch. Writing it in PyTorch is the point. Every step is differentiable, so we can ask
what change in a physical parameter would have changed the flood, and answer it with a gradient
rather than by trying values one at a time. Land cover classes map to Manning roughness, which is a
number describing how much a surface slows water down, and to infiltration, which is how fast water
soaks into the ground. Rainfall forcing comes from Open-Meteo reanalysis and forecast
\cite{openmeteo}.

\subsection{Metered infiltration, and a physics correction that made our floods worse}
An earlier version let every pervious cell soak up water at a constant rate for as long as the storm
lasted, which is wrong. Real soil fills up. Version~3 meters infiltration against a finite soil
storage capacity and caps the rate at the saturated hydraulic conductivity of the soil, which is
the fastest water can physically move through it. This change makes our predicted flooding
\emph{worse}, by design, and we committed the before and after comparison so the change is
auditable. At Patna and 100\,mm of rain, water that soaks in falls from 465\,449 to
391\,525\,m$^3$, ponded water rises by 1.4\,\%, and predicted flooded area goes from 20.63 to
20.78\,km$^2$. Soil fills to only 2.7\,\% of its capacity in that storm, which is the useful part
of the result. It says the earlier model was not merely mis-parameterised. It was letting a nearly
empty sponge absorb water forever.

\subsection{A millisecond emulator, and a reproducibility hazard}
A U-Net \cite{ronneberger2015} is trained to imitate the solver so the dashboard can answer what-if
questions instantly. Wet cells carry twenty times the weight in the loss, because flooding is rare
in space and a plain average error is minimised by predicting no flood anywhere. Validation error
across the bundles is 5 to 21\,cm and the dose-response is monotone, meaning more rain always
produces more water.

Building this exposed a failure worth reporting. Our dataset generator seeds the global random
number generator so that storm sampling is reproducible. That also makes the network's starting
weights bit-identical on every rebuild. On the Mumbai Island-City tile, where much of the window is
harbour, that particular starting point collapses deterministically. Within one epoch the optimizer
pushes the softplus output head so far negative that its gradient vanishes, and the emulator is
stuck predicting zero water everywhere. The validation error looks plausible at 16.3\,cm, which is
exactly the root mean square of the targets, and it stays frozen to the last digit for 40 epochs.
A frozen metric is the tell, because it means the gradients are exactly zero.

The important part is that rebuilding cannot fix this. We measured three bit-identical collapses
before diagnosing it. A fixed seed makes failures reproducible too. The shipped trainer now detects
a collapsed validation prediction, meaning no cell above the wetness threshold despite wet targets,
and retrains from an explicit alternate seed. The first alternate seed turned the worst tile into
the best-fitting emulator in the system, at 5.0\,cm validation error. We flag the pattern because
build pipelines that reseed globally for data reproducibility are common, and a frozen validation
metric is an easy and general warning sign.

\subsection{Serve layer, on free CPU hosting}
A FastAPI backend serves the committed per-city bundles, including cached OpenStreetMap road graphs
\cite{osm} and exposure counts, so the deployed service needs no external API access at run time.
It re-runs interventions live in seconds. A React and Leaflet front end exposes rainfall what-ifs,
ward alerts, drainage, storage and excavation planners, building and road level exposure, validation
panels, and a text brief grounded in the bundle's own numbers. A night lights layer
\cite{roman2018} overlays satellite-observed power outages. Everything from build to serve fits
inside free tiers, and a new city window takes roughly 15 to 25 minutes to build on a free T4 GPU,
including every artifact the server needs.

The system is deployed for 16 areas. These are Patna and two sub-windows, Bengaluru, six Mumbai
tiles, and six towns in Karnataka. A seventeenth area, Chennai, exists as a data bundle used for
the coastal experiment of Section~\ref{sec:coastal} but is not served, which is why the count of
deployed areas and the count of areas with radar ground truth are different numbers in this paper.
